\section{Introduction}\label{sec:intro}

%^C In one paragraph: introduce Software Supply Chain.
%^C How vital is it? (accelerate development; provide usable components)
%^C How severely has it been compromised?
%^C List examples of SSC attacks (xz-utils, log4j)
\par The Software Supply Chain (SSC) lies at the basis of modern digital
infrastructure. By providing software developers with reusable libraries,
packages, or components which can be assembled into a larger software
project\citep{sscdirection}, Software Supply Chain is integral to boosting
development efficacy, improving software quality, and reduce development cost
\citep{brooks1987no,mythicalmanmonth}. However, the number of Software Supply
Chain attacks, which come in the form of ransomware, backdoor, or vulnerability
injected into the source code, has grown exponentially over recent years.
This can be evidenced by the surging number of Common Vulnerabilities and
Exposures (CVEs), indicating the wide susceptibility of softwares. Also
Sonatype, a software supply chain platform, observed an annual increase of 156\%
in number of malicious packages, reaching a startling number of 704,102 identified
since 2019 \citep{sonatype2024report}. Such proliferation of malicious or
vulnerable components, coupled with exploding consumption of open-source 
packages and failure to promptly respond to reported vulnerabilities such as CVEs 
\citep{sonatype2024report}, have made Software Supply Chain
a broad attack surface, bringing about data or financial loss to
affected software developers and users.
For instance, the \texttt{xz-util} incident is an example of deliberately
injected backdoor, allowing the attacker to intercept sensitive data
or gain control of the affected machine \citep{xzincident}.
Despite all on-going efforts for security 
assurance and attack detection (e.g. \citep{konala25metadata,11022932,10695665}),
the sheer volume of risk to the
security of Software Supply Chain calls for new methods of governance.

%^C Introduce Large language model, in a more general way.
\par Recent advancements in Large Language Models (LLMs), including
improvement in model's intrinsic
capabilities\citep{llmsurvey,overviewllm,csllm} and the innovations of
enhancement techniques, contributed to their active involvement in
software engineering tasks, such as code suggestion, code completion, and
code review \citep{llm4se, llm4se2}. Because of their contextual reasoning
capacities and understanding of code intent, several studies have
explored leveraging these language models for guarding the security
of Software Supply Chain. For instance, by means of 
\textit{threat detection} (e.g. \citep{llmbackdoordet,singla2023empiricalstudyusinglarge,syed2025agenticaiautonomousdefense}),
malicious code such as backdoors in software packages can be 
reported; through enhanced \textit{security testing} and \textit{vulnerability detection}
(e.g. \citep{11022932,10695665,10628694}), vulnerabilities 
can be uncovered before exploitation; via automated \textit{vulnerability repair}
(e.g. \citep{vulrepairusingllm,vulmaster}), attackers have less time and chance
to exploit a reported vulnerability. LLM for Software Supply Chain
defense techniques will be discussed in Section~\ref{sec:rq2}.

\par Nevertheless, Large Language Models may be used equally for the good
and the evil. LLMs are susceptible to both intrinsic hallucinations and
a wide variety of attacks. Section~\ref{sec:rq3} discusses LLM's potential
of generating code with vague vulnerabilities or malice, because
of either internal hallucinations or deliberate attack on its
training data and prediction process. Furthermore, LLM's capability of
generating readable text has been leveraged for social engineering attack,
such as crafting phishing emails.

\par To the best of our knowledge, this study is the first to examine the
Large Language Model's implications of the security of Software Supply Chain.
It makes the following contributions:

\begin{itemize}
  \item A systematic literature review of LLM-based
  Software Supply Chain defense techniques (Section~\ref{sec:rq2}).

  \item A rigorous analysis of LLM's vulnerabilities and possible
  exploitations of LLM in Software Supply Chain attacks, in Section~\ref{sec:rq3}.

  \item Insights into mitigation of LLM-assisted Software Supply Chain
  attacks. They serve as guidance to responsible, ethical, and non-malignant
  use of LLM in Software Supply Chain.
\end{itemize}

